% 16-modules.tex — Clause 16: Module system
\chapter{Module system}
\label{sec:16-modules}
\index{module system}

\section{Overview}
\label{sec:mod.overview}

\pnum
Each Lain source file defines one module.  Modules are the unit of code
organization and compilation.  All names declared at the top level of a
source file are in that module's scope and are accessible from other
modules that import the module.

\section{Module identity}
\label{sec:mod.declaration}
\index{module declaration}

\pnum
A module's identity is \idbehav{derived from its file path by the
implementation}; there is \emph{no} \lain{module} declaration keyword (a
leading \lain{module}~\dots{} line is \diagcode{E100}).  The reference
implementation names a module by its path relative to the working directory,
with directory separators mapped to dots --- so \texttt{std/io.ln} is the
module \lain{std.io}.

\section{Import declarations}
\label{sec:mod.import}
\index{import declaration}

\pnum
\lain{import module.path} makes the declarations of the named module
available in the current source file.  The grammar production is
\gramref{import-declaration} (Annex~A).

\pnum
\textbf{Selective import.}  \lain{import module.path.\{name1, name2, ...\}}
brings only the listed names into scope rather than the whole module.  This is
the form used throughout the standard library, e.g.\
\lain{import std.mem.\{mem\_alloc, mem\_free\}} or
\lain{import std.fs.\{File, open\_file, close\_file\}}.

\pnum
The \lain{use} statement (\S\,\ref{sec:lex.keywords}) provides a
\emph{local-scope} import/alias inside a function body.  \idbehav{Its precise
semantics are implementation-defined in this revision}; the standard library
uses top-level \lain{import} rather than \lain{use}.

\pnum
\textbf{Path resolution}: a module path is converted to a file system path
by replacing dots with directory separators and appending \texttt{.ln}.
The search is relative to the compiler's working directory.
\idbehav{The exact search algorithm, including support for search paths,
is implementation-defined}.

\begin{longtable}{@{}ll@{}}
\toprule
\textbf{Module path} & \textbf{File path} \\
\midrule
\endfirsthead
\toprule
\textbf{Module path} & \textbf{File path} \\
\midrule
\endhead
\bottomrule
\endlastfoot
\lain{std.c}        & \texttt{std/c.ln}        \\
\lain{std.io}       & \texttt{std/io.ln}        \\
\lain{std.fs}       & \texttt{std/fs.ln}        \\
\lain{std.math}     & \texttt{std/math.ln}      \\
\lain{mylib.utils}  & \texttt{mylib/utils.ln}   \\
\end{longtable}

\section{Module loading}
\label{sec:mod.loading}
\index{module loading}

\pnum
When the compiler encounters an \lain{import} declaration:
\begin{enumerate}
  \item The module path is converted to a file path.
  \item If the module has already been loaded, the import is a no-op
        (deduplication).
  \item Otherwise, the file is read, lexed, parsed, and recursively
        resolved.
  \item The loaded module's declarations are merged into the importing
        module's declaration list.
\end{enumerate}

\pnum
Each module is loaded at most once.  Circular imports are therefore
implicitly handled: if module~A imports module~B which imports module~A,
the second import of~A is a no-op because A is already loaded.

\begin{note}
  The reference implementation maintains a linked list of loaded modules
  keyed by their dot-separated name.
\end{note}

\section{Name visibility}
\label{sec:mod.visibility}
\index{name visibility}

\pnum
By default, every name declared at module scope is \textit{public}:
visible to any module that imports the declaring module.  This default
matches the typical usage where module declarations form an API surface.

\pnum
A declaration prefixed with the attribute \tok{[private]}
(\S\,\ref{sec:10-declarations}) is \textit{module-private}: it is
visible only within the defining module and may not be referenced from
any importing module.

\begin{laincode}
// In module foo.bar:
func public_api() int { ... }    // public by default

[private]
func internal_helper() int { ... }   // only foo.bar may reference this
\end{laincode}

\begin{constraint}
  \textbf{Q-018.}
  A reference to a \tok{[private]} declaration from a module other than
  the one where it is defined is \illformed{}~(\diagcode{E084}).  The
  reference may appear in any expression position---identifier lookup,
  member access, type instantiation, etc.
\end{constraint}

\begin{note}
  Compilation under the reference implementation translates each
  \tok{[private]} declaration to a C declaration with internal linkage
  (\texttt{static} storage class), and each public declaration to one
  with external linkage.
\end{note}

\section{Standard library module paths}
\label{sec:mod.stdlib}
\index{standard library!module paths}

\pnum
The following module paths are defined by this specification for the
standard library:

\begin{longtable}{@{}ll@{}}
\toprule
\textbf{Module path} & \textbf{Description} \\
\midrule
\endfirsthead
\toprule
\textbf{Module path} & \textbf{Description} \\
\midrule
\endhead
\bottomrule
\endlastfoot
\lain{std.c}    & C standard library bindings \\
\lain{std.io}   & Console I/O \\
\lain{std.fs}   & File system operations \\
\lain{std.math} & Mathematical functions \\
\end{longtable}

\section{Name mangling}
\label{sec:mod.mangling}
\index{name mangling}

\pnum
The name of a Lain declaration in the translation output is
\idbehav{determined by the implementation}.  The reference implementation
uses the function name prefixed with the module path, with dots replaced
by underscores.  Name collisions between modules are
\idbehav{detected by the implementation; the reference implementation
appends a numeric suffix to resolve conflicts}.
